\section{Analisi critica}
\label{sec:discussion}

Questa sezione raccoglie le letture \emph{trasversali}, che nessun singolo esperimento
mostra da solo.

\subsection{SARSA vs Q-Learning: la teoria non predice il vincitore}
Il confronto ha prodotto esiti \emph{opposti} nei due ambienti, ed \`e istruttivo
proprio per questo. Su \textbf{CartPole} SARSA prevale su entrambi i fronti (return
$258$ vs $224$, deviazione standard $\pm 31$ vs $\pm 49$), mentre su \textbf{Acrobot}
Q\hbox{-}Learning ottiene un return leggermente migliore ($-341$ vs $-349$) ma SARSA
risulta quasi dieci volte pi\`u stabile tra i seed ($\pm 3.2$ vs $\pm 27.8$).

La differenza teorica \`e sempre la medesima: il target di Q\hbox{-}Learning impiega il
$\max$ e ignora l'esplorazione, quello di SARSA impiega l'azione realmente scelta e la
include. Le \emph{conseguenze}, tuttavia, dipendono dall'ambiente. L\`a dove
l'esplorazione \`e rischiosa (CartPole: una mossa casuale pu\`o far cadere l'asta),
SARSA guadagna perch\'e apprende una policy robusta al proprio rumore esplorativo. L\`a
dove il reward \`e sparso (Acrobot), il $\max$ di Q\hbox{-}Learning amplifica le rare
traiettorie fortunate che raggiungono l'obiettivo: quali siano dipende dal seed, il che
spiega i picchi pi\`u elevati ma anche la varianza molto superiore. Concludere che
``off-policy \`e meglio'' non ha pertanto valore predittivo: ci\`o che conta \`e come il
rischio dell'esplorazione interagisce con la struttura del reward.

\subsection{Perch\'e le curve non convergono mai davvero}
Nessuna curva tabellare si appiattisce del tutto, in lincea con le previsioni: con $\varepsilon$ e
$\alpha$ \emph{costanti} l'agente continua a esplorare e dunque la fluttuazione
asintotica \`e un comportamento atteso (S\&B \S2.5). Il fatto che l'annealing di
$\varepsilon$ riduca sensibilmente la varianza tra seed su entrambi gli ambienti ne
costituisce la conferma sperimentale (almeno funo a quanto il valore di $\varepsilon$ non si fissa nuovamente).

\subsection{Limiti metodologici riconosciuti}
\begin{itemize}
  \item \textbf{Predizione: misura locale.} Il confronto MC vs TD(0) traccia
        $V(s_{\text{ref}})$ di \emph{un} singolo stato; una misura globale (RMSE di $V$
        su tutti gli stati, come nel \emph{random walk} del corso) richiederebbe un $V$
        di riferimento (il ground truth) che in questi ambienti non \`e disponibile a causa della struttura
        del problema, con un numero finito di stati da valutare.
  \item \textbf{Asse x in ``episodi'' nell'ablazione Pong.} A parit\`a di 200k step le
        diverse configurazioni producono numeri diversi di episodi, per cui confrontarle
        allo stesso indice non costituisce un confronto a perfetta parit\`a di calcolo. Un
        asse in step ambiente sarebbe pi\`u rigoroso; la conclusione qualitativa non
        cambia, dal momento che le curve senza replay restano piatte per l'intera durata.
  \item \textbf{Bande assenti nelle figure di sensibilit\`a}, per ragioni di
        leggibilit\`a (Sez.~\ref{sec:experiments-reading}).
  \item \textbf{Il DQN \`e valutato durante il training}, con $\varepsilon > 0$: si
        tratta di una stima leggermente \emph{pessimistica} della policy greedy. I
        checkpoint sono salvati, e una valutazione greedy separata sarebbe fattibile senza
        necessit\`a di riaddestrare.
\end{itemize}

\subsection{Problemi implementativi incontrati}
Distinti dai limiti intrinseci del RL, questi sono ostacoli concreti che \`e stato
necessario risolvere:
\begin{itemize}
  \item \textbf{GPU Blackwell (RTX 5060)}: richiede la build PyTorch CUDA 12.8; con
        build precedenti la GPU non viene rilevata e il training ricade silenziosamente
        su CPU.
  \item \textbf{Pipeline Atari, tre insidie}: (i) l'ambiente \code{-v5} possiede gi\`a
        un frame-skip interno che, combinato con \code{AtariPreprocessing}, produrrebbe
        uno skip effettivo di 16 frame, rendendo necessario forzare
        \code{frameskip=1}; (ii) in Gymnasium 1.x \code{FrameStack} si chiama
        \code{FrameStackObservation}; (iii) i frame impilati sono gi\`a in formato
        channel-first \code{(4,84,84)} e non devono essere trasposti. Un errore su uno
        qualsiasi dei tre produce un crash oppure, peggio ancora, input silenziosamente
        malformati.
  \item \textbf{Memoria del replay buffer}: 1M di transizioni Atari
        richiederebbero $\approx 28$\,GB, pertanto il buffer \`e stato limitato a 100k
        transizioni mantenendo i frame in formato \code{uint8}.
  \item \textbf{DQN non riproducibile per un RNG nascosto.} Run che avrebbero dovuto
        essere identiche divergevano: l'azione esplorativa viene estratta da
        \code{env.action\_space.sample()}, e lo spazio delle azioni dispone di un
        \emph{RNG separato} che n\'e \code{set\_seed} n\'e \code{reset(seed=)}
        inizializzano. Nelle fasi iniziali ($\varepsilon\!\approx\!1$) quasi tutta
        l'esperienza proveniva da quel RNG non seminato. Aggiunto
        \code{env.action\_space.seed(seed)}; i risultati Pong riportati provengono da
        questo run corretto e riproducibile.
  \item \textbf{Episodio troncato contato come completo.} In DQN, quando il budget di step si
        esauriva a met\`a partita, il return parziale (troppo ottimistico) veniva
        registrato come episodio completo, producendo un picco spurio verso l'alto sul
        bordo destro di ogni curva. Corretto scartando l'episodio non terminato.
\end{itemize}
